\documentclass[11pt]{article}
\usepackage[margin=1in]{geometry}
\usepackage{amsmath,amssymb,amsthm,mathtools}
\usepackage{enumitem}
\usepackage{booktabs}
\usepackage{hyperref}

\newtheorem{theorem}{Theorem}[section]
\newtheorem{lemma}[theorem]{Lemma}
\newtheorem{proposition}[theorem]{Proposition}
\newtheorem{definition}[theorem]{Definition}
\newtheorem{remark}[theorem]{Remark}
\newtheorem{corollary}[theorem]{Corollary}

\newcommand{\K}{\mathsf K}
\newcommand{\A}{\mathsf A}
\newcommand{\E}{\mathsf E}
\newcommand{\Y}{\mathcal Y}
\newcommand{\Hil}{\mathcal H}
\newcommand{\Fix}{\operatorname{Fix}}
\newcommand{\Ran}{\operatorname{Ran}}
\newcommand{\tr}{\operatorname{tr}}
\newcommand{\pre}{\preceq}
\newcommand{\suc}{\succeq}

\title{Step 34: Completed Explicit-Formula Bridge Decomposition}
\author{Anti-Localization Working Notes}
\date{May 2026}

\begin{document}
\maketitle

\begin{abstract}
Step 33 isolated the RH-facing bridge as a contractive domination
\[
  \A_Z \preceq \K^-.
\]
Equivalently, if \(V_Z:\Y^-\to Z\) is the anti-invariant zero-displacement readout and \(W_E:\Y^-\to E\) is the completed carrier-side minimal-energy readout, then the exact domination square is the existence of a contraction \(T:E\to Z\) such that
\[
  V_Z=T W_E.
\]
This note decomposes that single bridge into the visible pieces that an explicit formula must supply: prime-power, archimedean/gamma, pole/completion, and support-tail terms. The purpose is not to prove RH. The purpose is to state exactly what the completed explicit-formula bridge must certify, how defects enter, and why trace or invariant-sector agreement is insufficient.
\end{abstract}

\section{Setup}

Let \(J\) be the completed anti-linear duality. After realification of the response space, write
\[
  \Y_{\mathbb R}=\Y^+\oplus \Y^-,
  \qquad
  P_\pm=\frac{I\pm J}{2}.
\]
For RH, \(J(s)=1-\overline{s}\), and \(\Fix(J)=\{\Re(s)=1/2\}\). The anti-invariant space \(\Y^-\) is the response space in which off-critical displacement is visible.

Let
\[
  V_Z:\Y^-\to \mathcal Z
\]
be the zero-side anti-invariant displacement readout. Its zero-side displacement matrix is
\[
  \A_Z:=V_Z^*V_Z.
\]
Let the completed carrier-side anti-invariant currency be represented by a feature map
\[
  W_E:\Y^-\to \mathcal E,
  \qquad
  \K^-:=W_E^*W_E.
\]
Step 33 established the basic domination/factorization equivalence:
\[
  \A_Z\preceq \K^-
  \quad\Longleftrightarrow\quad
  \exists T:\mathcal E\to\mathcal Z,\ \|T\|\le 1,
  \quad V_Z=T W_E.
\]
The present note asks what it means for the completed explicit formula to provide \(W_E\) in pieces.

\section{Completed bridge components}

\begin{definition}[Completed explicit-formula bridge package]
A completed explicit-formula bridge package on \(\Y^-\) consists of Hilbert spaces
\[
  \mathcal E_{\rm pr},\quad
  \mathcal E_{\infty},\quad
  \mathcal E_{\rm pole},\quad
  \mathcal E_{\rm tail},
\]
and feature maps
\[
  W_{\rm pr}:\Y^-\to\mathcal E_{\rm pr},
  \quad
  W_\infty:\Y^-\to\mathcal E_\infty,
  \quad
  W_{\rm pole}:\Y^-\to\mathcal E_{\rm pole},
  \quad
  W_{\rm tail}:\Y^-\to\mathcal E_{\rm tail}.
\]
The completed feature map is the column operator
\[
  W_{\rm cmp}
  :=
  \begin{bmatrix}
    W_{\rm pr}\\ W_\infty\\ W_{\rm pole}\\ W_{\rm tail}
  \end{bmatrix}
  :\Y^-\to
  \mathcal E_{\rm pr}\oplus\mathcal E_\infty\oplus\mathcal E_{\rm pole}\oplus\mathcal E_{\rm tail},
\]
with completed currency
\[
  \K_{\rm cmp}
  :=W_{\rm cmp}^*W_{\rm cmp}
  =\K_{\rm pr}+\K_\infty+\K_{\rm pole}+\K_{\rm tail},
\]
where \(\K_\bullet=W_\bullet^*W_\bullet\).
\end{definition}

\begin{remark}[Raw terms versus completed features]
The explicit formula is often written as a signed balance of zero, prime, gamma, and pole terms. A signed balance is not yet a currency certificate. A currency certificate requires feature maps \(W_\bullet\) or an equivalent positive quadratic representation. Individual raw prime/gamma/pole expressions need not be positive. The completed carrier must assemble them into a positive anti-invariant currency \(\K_{\rm cmp}\).
\end{remark}

\section{Contractive assembly theorem}

\begin{theorem}[Completed bridge decomposition]
Let \(V_Z:\Y^-\to\mathcal Z\) and let \(W_{\rm cmp}\) be the completed column feature map. The following are equivalent:
\begin{enumerate}[label=(\roman*)]
  \item \(\A_Z\preceq \K_{\rm cmp}\).
  \item There exists a contraction
  \[
    T:\mathcal E_{\rm pr}\oplus\mathcal E_\infty\oplus\mathcal E_{\rm pole}\oplus\mathcal E_{\rm tail}\to \mathcal Z
  \]
  such that
  \[
    V_Z=T W_{\rm cmp}.
  \]
  \item There are component maps \(T_\bullet:\mathcal E_\bullet\to\mathcal Z\) forming a row contraction
  \[
    T=[T_{\rm pr}\ T_\infty\ T_{\rm pole}\ T_{\rm tail}],
    \qquad \|T\|\le1,
  \]
  such that
  \[
    V_Z=T_{\rm pr}W_{\rm pr}+T_\infty W_\infty+T_{\rm pole}W_{\rm pole}+T_{\rm tail}W_{\rm tail}.
  \]
\end{enumerate}
\end{theorem}

\begin{proof}
The equivalence of (ii) and (iii) is just the row/column decomposition of a bounded operator on a direct sum. The equivalence of (i) and (ii) is Douglas factorization: for bounded operators \(V_Z\) and \(W_{\rm cmp}\),
\[
  V_Z^*V_Z\preceq W_{\rm cmp}^*W_{\rm cmp}
\]
if and only if \(V_Z=T W_{\rm cmp}\) for some contraction \(T\). In finite dimensions this follows immediately from the range-inclusion/pseudoinverse construction \(T=V_ZW_{\rm cmp}^\dagger\) on \(\overline{\Ran W_{\rm cmp}}\) and extension by zero.
\end{proof}

\begin{corollary}[Zero confinement with completed bridge]
If the completed explicit-formula bridge gives \(\A_Z\preceq\K_{\rm cmp}\) and the accepted anti-invariant budget gives
\[
  \K_{\rm cmp}\preceq\Theta^-,
\]
then
\[
  \A_Z\preceq\Theta^-.
\]
If \(\Theta^-=0\), all visible zero mass lies in \(\Fix(J)\).
\end{corollary}

\section{Defective bridge decomposition}

The honest finite or incomplete bridge is not usually exact. Missing completion terms, tails, illegal support, endpoint losses, and protocol mismatches must be represented as defects.

\begin{definition}[Bridge defect]
A defective completed bridge consists of \(W_{\rm cmp}\), a row contraction \(T\), and a residual readout \(R:\Y^-\to\mathcal Z\) such that
\[
  V_Z=T W_{\rm cmp}+R.
\]
If
\[
  R^*R\preceq E_{\rm EF},
\]
then \(E_{\rm EF}\) is the explicit-formula bridge defect.
\end{definition}

\begin{theorem}[Defective bridge bound]
If
\[
  V_Z=T W_{\rm cmp}+R,
  \qquad \|T\|\le1,
  \qquad R^*R\preceq E_{\rm EF},
\]
then for every \(t>0\),
\[
  \A_Z\preceq (1+t)\K_{\rm cmp}+(1+t^{-1})E_{\rm EF}.
\]
In particular, if the exact cross term is recorded or orthogonalized so that
\[
  (TW_{\rm cmp})^*R+R^*(TW_{\rm cmp})\preceq0,
\]
then the sharper bound
\[
  \A_Z\preceq \K_{\rm cmp}+E_{\rm EF}
\]
holds.
\end{theorem}

\begin{proof}
For each \(y\in\Y^-\),
\[
  \|V_Zy\|^2=\|TW_{\rm cmp}y+Ry\|^2
  \le (1+t)\|TW_{\rm cmp}y\|^2+(1+t^{-1})\|Ry\|^2.
\]
Since \(T\) is a contraction,
\[
  \|TW_{\rm cmp}y\|^2\le \|W_{\rm cmp}y\|^2=y^*\K_{\rm cmp}y.
\]
Using \(R^*R\preceq E_{\rm EF}\) gives the first claim. The sharper claim follows by expanding the square and using the assumed cross-term record.
\end{proof}

\begin{corollary}[Defective zero confinement]
If
\[
  \K_{\rm cmp}\preceq\Theta^-,
\]
then every defective bridge as above gives
\[
  \A_Z\preceq (1+t)\Theta^-+(1+t^{-1})E_{\rm EF}.
\]
Thus missing completion terms appear as quantitative off-fixed-locus allowance.
\end{corollary}

\section{Component obligations}

The bridge package is accepted only when each component is visible, lawful, and has a defect record.

\subsection{Prime-power component}
The prime-power feature \(W_{\rm pr}\) carries the lower prime trace layer. It must record:
\begin{enumerate}[label=\textup{Pr\arabic*.}]
  \item the prime-power weights, including \(\Lambda(n)\), \(\log n\), and support/window factors;
  \item the threshold rule for when a prime power enters the support ledger;
  \item the staging rule matching finite prime sums with the test-function support;
  \item the route compatibility record between additive and multiplicative descriptions;
  \item the tail defect \(E_{\rm pr,tail}\) for omitted prime powers.
\end{enumerate}
The prime component alone is not a completed currency unless it is positive on \(\Y^-\) and visibly anti-invariant. Otherwise it is support evidence.

\subsection{Archimedean/gamma component}
The archimedean feature \(W_\infty\) carries the gamma and real-place completion. It must record:
\begin{enumerate}[label=\textup{Ga\arabic*.}]
  \item the completed gamma contribution, not a raw zeta-only shadow;
  \item its anti-linear duality behavior under \(J(s)=1-\overline{s}\);
  \item its action near legal zero modes and slow modes;
  \item its defect if approximated by finite quadrature or asymptotic expansion.
\end{enumerate}
Dropping \(W_\infty\) is a completion defect, not a harmless normalization change.

\subsection{Pole/completion component}
The pole/completion feature \(W_{\rm pole}\) carries the terms required by \(s(s-1)\), endpoint residues, legal null modes, and quotienting. It must record:
\begin{enumerate}[label=\textup{Po\arabic*.}]
  \item the legal null-mode quotient;
  \item annihilation of probes on the legal kernel;
  \item pole separation from nontrivial zero displacement;
  \item endpoint or boundary defects if support is truncated.
\end{enumerate}
Without this component, the anti-invariant displacement can couple to legal zero-energy directions and create infinite capacity.

\subsection{Support/tail component}
The support/tail feature \(W_{\rm tail}\) records finite-support, truncation, smoothing, and finite-to-continuum defects:
\begin{enumerate}[label=\textup{Ta\arabic*.}]
  \item test-function admissibility and support closure;
  \item endpoint smoothing losses;
  \item finite zero/prime cutoff errors;
  \item response-transport defects between finite and exact carriers;
  \item the dominated or summable envelope needed for promotion.
\end{enumerate}
This component is the bridge between a finite calculation and an exact claim.

\section{Trace-shadow and invariant-sector warnings}

\begin{proposition}[Trace agreement is not domination]
The scalar identity
\[
  \tr \A_Z=\tr \K_{\rm cmp}
\]
does not imply
\[
  \A_Z\preceq\K_{\rm cmp}.
\]
\end{proposition}

\begin{proof}
Take
\[
  \A_Z=\begin{pmatrix}2&0\\0&0\end{pmatrix},
  \qquad
  \K_{\rm cmp}=\begin{pmatrix}1&0\\0&1\end{pmatrix}.
\]
Both traces are \(2\), but \(\A_Z\npreceq\K_{\rm cmp}\).
\end{proof}

\begin{proposition}[Invariant sector is not anti-invariant control]
Control of the invariant sector \(\Y^+\) does not imply control of the anti-invariant sector \(\Y^-\).
\end{proposition}

\begin{proof}
On \(\Y^+\oplus\Y^-\), take
\[
  \K=\begin{pmatrix}1/2&0\\0&2\end{pmatrix},
  \qquad
  \Theta=I.
\]
The invariant sector passes, but the anti-invariant sector fails.
\end{proof}

\section{Square-ledger statuses}

A completed explicit-formula bridge has status \texttt{accepted} only if:
\begin{enumerate}[label=\textup{G\arabic*.}]
  \item the zero-side readout \(V_Z\) is declared on \(\Y^-\);
  \item all four feature components \(W_{\rm pr},W_\infty,W_{\rm pole},W_{\rm tail}\) are declared or marked inapplicable with record;
  \item the assembled feature map \(W_{\rm cmp}\) is a lawful carrier-side witness, not a public shadow;
  \item the contractive factorization \(V_Z=T W_{\rm cmp}\) is proved, or the defect \(E_{\rm EF}\) is recorded;
  \item the anti-invariant budget \(\K_{\rm cmp}\preceq\Theta^-\) is accepted by the anti-localization test-family gate;
  \item finite truncation and support promotion have a probe-wise/predictive defect record;
  \item no trace-only, invariant-only, or raw-prime-only statement is overread as the completed bridge.
\end{enumerate}

The status is \texttt{support-only} if the pieces suggest the bridge but no contraction or defect is proved. It is \texttt{overread} if trace agreement or invariant-sector agreement is claimed as anti-invariant domination. It is \texttt{failed-completion} if gamma, pole, support, or tail records are missing from a claim that requires them.

\section{RH-facing statement}

For RH, the completed explicit-formula bridge must provide:
\[
  V_Z=T
  \begin{bmatrix}
    W_{\rm pr}\\ W_\infty\\ W_{\rm pole}\\ W_{\rm tail}
  \end{bmatrix},
  \qquad \|T\|\le1,
\]
on the anti-invariant response space for \(J(s)=1-\overline{s}\). If the accepted anti-localization test family gives
\[
  \K_{\rm pr}+\K_\infty+\K_{\rm pole}+\K_{\rm tail}\preceq\Theta^-,
\]
then
\[
  \A_Z\preceq\Theta^-.
\]
If \(\Theta^-=0\), the visible zero ledger is confined to the critical line. If defects are present,
\[
  \A_Z\preceq (1+t)\Theta^-+(1+t^{-1})E_{\rm EF},
\]
so the defect is exactly the remaining off-critical allowance.

\begin{remark}[Nonclaim]
This schema does not prove RH. It states what the completed explicit formula must supply for the zero-confinement theorem to apply. In particular, it does not license raw prime truncations, trace identities, or invariant-sector passivity as anti-invariant zero confinement.
\end{remark}

\end{document}
